\documentclass[12pt,a4paper]{article}

% ============== PACKAGES ==============
\usepackage[utf8]{inputenc}
\usepackage[T1]{fontenc}
\usepackage{amsmath,amssymb,amsthm}
\usepackage{physics}
\usepackage{graphicx}
\usepackage{booktabs}
\usepackage{array}
\usepackage{tabularx}
\usepackage{geometry}
\usepackage{hyperref}
\usepackage{cleveref}
\usepackage{xcolor}
\usepackage{tcolorbox}
\usepackage{fancyhdr}
\usepackage{titlesec}
\usepackage{enumitem}
\usepackage{float}

% ============== PAGE SETUP ==============
\geometry{margin=2.5cm}
\hypersetup{
    colorlinks=true,
    linkcolor=blue!70!black,
    citecolor=green!50!black,
    urlcolor=blue!70!black
}

% ============== CUSTOM BOXES ==============
\tcbuselibrary{skins,breakable}

\newtcolorbox{keyresult}{
    colback=blue!5!white,
    colframe=blue!75!black,
    fonttitle=\bfseries,
    title=Key Result,
    breakable
}

\newtcolorbox{insight}{
    colback=green!5!white,
    colframe=green!60!black,
    fonttitle=\bfseries,
    title=Physical Insight,
    breakable
}

\newtcolorbox{prediction}{
    colback=orange!5!white,
    colframe=orange!70!black,
    fonttitle=\bfseries,
    title=Testable Prediction,
    breakable
}

\newtcolorbox{dedication}{
    colback=yellow!10!white,
    colframe=yellow!60!black,
    fonttitle=\bfseries,
    title=Dedication to Humanity,
    breakable
}

% ============== HEADER/FOOTER ==============
\pagestyle{fancy}
\fancyhf{}
\fancyhead[L]{\small Coherent Bulk Focusing}
\fancyhead[R]{\small EDC Fusion Paper}
\fancyfoot[C]{\thepage}
\renewcommand{\headrulewidth}{0.4pt}

% ============== TITLE ==============
\title{
    \vspace{-1cm}
    {\LARGE\textbf{Coherent Bulk Focusing: A 5D Approach to Engineered Nuclear Fusion}}\\[0.5cm]
    {\large From Quantum Tunneling to Controlled Energy Production}
}
\author{
    \textbf{Igor Grčman}\\[0.2cm]
    \small Independent Researcher\\
    \small \href{mailto:igor@edctheory.org}{igor@edctheory.org}
}
\date{January 2026}

% ============== DOCUMENT ==============
\begin{document}

\maketitle

\begin{dedication}
\textbf{This work is dedicated to humanity.} The knowledge contained herein is released freely and may not be patented, restricted, or monopolized. Energy belongs to everyone.
\end{dedication}

\begin{abstract}
We propose a revolutionary approach to nuclear fusion based on the five-dimensional geometry of Elastic Diffusive Cosmology (EDC). Since quantum tunneling represents particle transit through the 5D Bulk rather than mysterious barrier penetration, we can engineer conditions that facilitate this transit. By creating coherent wave superposition directed into the Bulk---analogous to phased-array focusing in 3D acoustics---we can establish focal points where the membrane-Bulk barrier is minimized. Particles at these focal points experience dramatically enhanced tunneling probability, enabling fusion at temperatures far below those required by brute-force thermal approaches. This paper presents the theoretical framework, proposes specific experimental configurations, and outlines a path toward clean, abundant energy for humanity.

\vspace{0.3cm}
\noindent\textbf{Keywords:} Nuclear fusion, quantum tunneling, extra dimensions, phased array, coherent focusing, clean energy, EDC
\end{abstract}

\tableofcontents
\newpage

%==============================================================================
\section{Introduction}
%==============================================================================

\subsection{The Current State of Fusion}

After 70 years and billions of dollars, controlled nuclear fusion remains ``30 years away.'' The fundamental approach has not changed:

\begin{enumerate}
    \item Heat plasma to extreme temperatures ($\sim$150 million K)
    \item Confine it magnetically or inertially
    \item Hope enough particles randomly tunnel through the Coulomb barrier
    \item Extract net energy
\end{enumerate}

This brute-force approach treats tunneling as an immutable quantum probability. We heat particles because we know no other way to increase tunneling rates.

\subsection{The EDC Paradigm Shift}

Elastic Diffusive Cosmology reveals that quantum tunneling is not mysterious wavefunction behavior but \textbf{geometric transit through a higher dimension}:

\begin{itemize}
    \item The Coulomb barrier exists on the 3D membrane
    \item The 5D Bulk contains no electrostatic barrier
    \item Particles can transit through the Bulk to bypass barriers
    \item \textbf{Tunneling probability depends on Bulk accessibility}
\end{itemize}

This transforms fusion from a probability game to an \textbf{engineering challenge}: How do we facilitate Bulk transit?

\subsection{The Key Insight}

In 3D, we routinely use wave superposition to create focused effects:
\begin{itemize}
    \item Ultrasonic levitation (acoustic focusing)
    \item Phased-array radar (electromagnetic focusing)
    \item HIFU therapy (ultrasound surgery)
\end{itemize}

\begin{insight}
\textbf{What if we could focus waves into the Bulk?}

Coherent sources on the membrane, properly phased, could create focal points in the 5D Bulk where transit is facilitated. Particles near these focal points would experience enhanced tunneling---\textbf{engineered fusion}.
\end{insight}

%==============================================================================
\section{Theoretical Foundation}
%==============================================================================

\subsection{The 5D Geometry}

Our universe is a 3D membrane ($\Sigma$) embedded in a 5D Bulk:
\begin{equation}
X^A = (w, x, y, z, \xi)
\end{equation}

\begin{itemize}
    \item $(x, y, z)$: Our spatial dimensions
    \item $w$: Depth into Bulk (perpendicular to membrane)
    \item $\xi$: Compact dimension (radius $R_\xi \approx 2.16 \times 10^{-18}$ m)
\end{itemize}

The membrane sits at $w = 0$. Particles are topological defects bound to the membrane.

\subsection{Membrane Dynamics}

The membrane has:
\begin{itemize}
    \item \textbf{Tension:} $\sigma_{\text{eff}} = 1.41 \times 10^{18}$ J/m$^2$
    \item \textbf{Thickness:} $R_\xi \approx 2.16 \times 10^{-18}$ m
    \item \textbf{Vibrational modes:} Waves can propagate on and through the membrane
\end{itemize}

The membrane is not rigid---it vibrates, deforms, and interacts with the Bulk (Plenum).

\subsection{Wave Propagation}

Waves on the membrane satisfy:
\begin{equation}
\left(\nabla^2 - \frac{1}{c_m^2}\frac{\partial^2}{\partial t^2}\right)\phi = 0
\end{equation}

where $c_m$ is the membrane wave speed.

Critically, membrane vibrations \textbf{extend into the Bulk}. A vibrating membrane displaces the surrounding Plenum, creating waves in the $w$-direction.

\subsection{Characteristic Frequencies}

The fundamental frequency associated with membrane thickness:
\begin{equation}
f_0 = \frac{c}{R_\xi} \approx \frac{3 \times 10^8}{2 \times 10^{-18}} \approx 1.4 \times 10^{26} \text{ Hz}
\end{equation}

This is gamma-ray frequency---impractical for direct excitation.

However, \textbf{subharmonics} exist:
\begin{equation}
f_n = \frac{f_0}{n}
\end{equation}

\begin{table}[H]
\centering
\begin{tabular}{@{}lll@{}}
\toprule
\textbf{n} & \textbf{Frequency} & \textbf{Regime} \\
\midrule
$10^{20}$ & $10^6$ Hz & MHz (radio) \\
$10^{18}$ & $10^8$ Hz & 100 MHz (RF) \\
$10^{14}$ & $10^{12}$ Hz & THz (far-IR) \\
$10^{10}$ & $10^{16}$ Hz & UV \\
\bottomrule
\end{tabular}
\caption{Subharmonic frequencies potentially coupling to membrane modes.}
\end{table}

Accessible frequencies may couple to membrane modes through nonlinear resonance.

%==============================================================================
\section{Coherent Bulk Focusing}
%==============================================================================

\subsection{The Concept}

In 3D phased arrays, multiple coherent sources create interference patterns. By controlling the phase of each source, we can focus energy at arbitrary points.

\textbf{We propose the same principle in 5D:}

Multiple coherent excitations on the membrane, properly phased, create a focal point \textbf{in the Bulk} (at some depth $w_f > 0$).

\begin{figure}[H]
\centering
\begin{verbatim}
PHASED ARRAY IN 3D (conventional):

    Source 1 ──→ )))
    Source 2 ──→ )))  → FOCAL POINT in 3D
    Source 3 ──→ )))


PHASED ARRAY INTO BULK (EDC):

    Membrane (w = 0)
    ════●════●════●════●════●════
        S₁   S₂   S₃   S₄   S₅
         ↘    ↘    ↓    ↙    ↙
          ↘    ↘   ↓   ↙    ↙
           ↘    ↘  ↓  ↙    ↙
            ↘    ↘ ↓ ↙    ↙
             ↘    ●    ↙   ← FOCAL POINT in Bulk
              ↘      ↙       (w = w_f)
    Bulk (w > 0)
\end{verbatim}
\caption{Comparison of conventional 3D phased array with EDC Bulk focusing.}
\end{figure}

\subsection{Mathematical Description}

For $N$ sources at positions $\vec{r}_i$ on the membrane, each emitting with amplitude $A_i$ and phase $\phi_i$:
\begin{equation}
\Psi(\vec{r}, w, t) = \sum_{i=1}^{N} A_i \frac{\exp(ik|\vec{R}_i|)}{|\vec{R}_i|} \cos(\omega t + \phi_i)
\end{equation}

where $\vec{R}_i = (\vec{r} - \vec{r}_i, w)$ is the 4D displacement to the field point.

At the focal point $(\vec{r}_f, w_f)$, phases are chosen so all contributions add constructively:
\begin{equation}
\phi_i = -k|\vec{R}_{i,f}| + \phi_0
\end{equation}

where $\vec{R}_{i,f}$ is the displacement from source $i$ to the focal point.

\subsection{Intensity at Focus}

With perfect phasing, the intensity at focus scales as:
\begin{equation}
I_f \propto N^2 A^2
\end{equation}

The $N^2$ scaling (vs. $N$ for incoherent sources) is the hallmark of coherent focusing.

\subsection{Effect on Tunneling}

At the focal point, the membrane-Bulk interface is \textbf{locally excited}. This reduces the effective barrier for particle detachment from the membrane.

The modified tunneling probability:
\begin{equation}
P_{\text{tunnel}}^{(\text{focused})} = P_{\text{tunnel}}^{(0)} \times \exp\left(\frac{\Delta E_{\text{focus}}}{k_B T_{\text{eff}}}\right)
\end{equation}

where $\Delta E_{\text{focus}}$ is the energy contributed by the focused field.

\begin{keyresult}
Even modest focusing can dramatically enhance tunneling due to the exponential dependence.
\end{keyresult}

%==============================================================================
\section{Physical Implementation}
%==============================================================================

\subsection{Candidate Excitation Mechanisms}

What physical processes can excite membrane modes?

\textbf{Mechanism A: Piezoelectric Transducers}
\begin{itemize}
    \item Piezoelectric crystals convert electrical signals to mechanical vibrations
    \item Frequencies: kHz to GHz achievable
    \item Arrays: Well-developed phased-array technology
    \item Control: Precise phase and amplitude control
\end{itemize}

\textbf{Mechanism B: Electromagnetic Fields}
\begin{itemize}
    \item Oscillating EM fields couple to charged particles on the membrane
    \item High frequencies accessible (THz with current technology)
    \item Can be focused independently
    \item May couple to membrane through charge motion
\end{itemize}

\textbf{Mechanism C: Laser Pulses}
\begin{itemize}
    \item Ultrashort laser pulses create impulsive excitations
    \item Femtosecond pulses = broad frequency spectrum
    \item Coherent control possible
    \item High peak intensities
\end{itemize}

\textbf{Mechanism D: Plasma Oscillations}
\begin{itemize}
    \item Collective modes in plasma (Langmuir waves, ion acoustic waves)
    \item Natural coupling to charged particles
    \item Self-consistent with fusion environment
    \item Frequencies depend on plasma parameters
\end{itemize}

\subsection{Proposed Experimental Configuration}

\begin{figure}[H]
\centering
\begin{verbatim}
EXPERIMENTAL SETUP (Top View):

                    ┌─────────────────┐
                    │   VACUUM CHAMBER │
                    │                  │
         P ────●    │      ┌───┐      │    ● ──── P
              ╱     │      │ T │      │     ╲
             ╱      │      │ A │      │      ╲
        P ──●       │      │ R │      │       ●── P
            │       │      │ G │      │       │
            │       │      │ E │      │       │
        P ──●       │      │ T │      │       ●── P
             ╲      │      │   │      │      ╱
              ╲     │      └───┘      │     ╱
         P ────●    │                  │    ● ──── P
                    │                  │
                    └─────────────────┘
                    
    P = Piezoelectric transducers (phased array)
    Target = Deuterium/Tritium pellet or gas
\end{verbatim}
\caption{Proposed experimental setup for coherent Bulk focusing.}
\end{figure}

\begin{figure}[H]
\centering
\begin{verbatim}
EXPERIMENTAL SETUP (Side View):

    ════════════════════════════════════════
              Piezo array (top)
    ────────────────────────────────────────
                    ↓ ↓ ↓ ↓ ↓
                     ↓ ↓ ↓
                      ↓ ↓
                       ↓
                    [TARGET]  ← Focal point
                       ↑       (in Bulk sense)
                      ↑ ↑
                     ↑ ↑ ↑
                    ↑ ↑ ↑ ↑ ↑
    ────────────────────────────────────────
              Piezo array (bottom)
    ════════════════════════════════════════
\end{verbatim}
\caption{Side view showing wave focusing toward target.}
\end{figure}

\subsection{Experimental Parameters}

\begin{table}[H]
\centering
\begin{tabular}{@{}lll@{}}
\toprule
\textbf{Parameter} & \textbf{Value} & \textbf{Rationale} \\
\midrule
Array geometry & Spherical & Optimal 3D focusing \\
Number of elements & 100--1000 & Sufficient for sharp focus \\
Frequency & Scan 1 kHz -- 1 GHz & Search for resonances \\
Phase control & $<1°$ precision & Coherent addition \\
Target & D$_2$ or DT & Standard fusion fuel \\
Diagnostics & Neutron, $\gamma$, heat & Fusion signatures \\
\bottomrule
\end{tabular}
\caption{Experimental parameters for initial tests.}
\end{table}

\subsection{Control Variables}

To test the EDC hypothesis, systematically vary:

\begin{enumerate}
    \item \textbf{Frequency} --- Scan for resonant enhancement
    \item \textbf{Phase configuration} --- Compare focused vs. random
    \item \textbf{Number of sources} --- Test $N^2$ scaling
    \item \textbf{Geometry} --- Different focal depths
    \item \textbf{Target material} --- Different nuclear reactions
\end{enumerate}

%==============================================================================
\section{Resonance Conditions}
%==============================================================================

\subsection{Finding the Right Frequency}

The membrane has characteristic frequencies. We don't know them precisely, but we can search.

\textbf{Signature of resonance:}
\begin{itemize}
    \item Sharp increase in fusion rate at specific frequency
    \item Q-factor indicating resonant width
    \item Harmonic structure (multiples of fundamental)
\end{itemize}

\subsection{Predicted Resonance Structure}

If the membrane has fundamental mode $f_0$, we expect resonances at:
\begin{equation}
f_n = \frac{f_0}{n}, \quad f_m = m \cdot f_{\text{sub}}
\end{equation}

where $f_{\text{sub}}$ is a subharmonic that couples to accessible frequencies.

\subsection{Scanning Protocol}

\begin{verbatim}
FREQUENCY SCAN PROTOCOL:

1. Start at f = 1 kHz
2. Measure fusion rate (neutron count)
3. Increment frequency by Δf
4. Repeat until f = 1 GHz
5. Identify peaks
6. Fine-scan around peaks
7. Map resonance structure

    Fusion
    Rate
      ↑
      │           ╱╲
      │          ╱  ╲
      │    ╱╲   ╱    ╲
      │   ╱  ╲ ╱      ╲
      │──╱────╳────────╲────────→ Frequency
      │ ╱              ╲
      │╱                ╲
      └─────────────────────────→
         f₁    f₂    f₃
         
    Resonant peaks indicate membrane modes
\end{verbatim}

%==============================================================================
\section{Alternative Configurations}
%==============================================================================

\subsection{Crystal Lattice Focusing}

Certain crystal structures may naturally create Bulk focusing:

\textbf{Palladium hydride:}
\begin{itemize}
    \item FCC lattice with D in octahedral sites
    \item Regular geometry may create standing wave pattern
    \item Historical anomalies (Fleischmann-Pons) may be explained
\end{itemize}

\textbf{Lithium tantalate (LiTaO$_3$):}
\begin{itemize}
    \item Pyroelectric: generates E-fields with temperature change
    \item Already demonstrated anomalous neutron production
    \item Crystal structure may facilitate focusing
\end{itemize}

\subsection{Acoustic Cavitation}

Ultrasonic cavitation creates imploding bubbles:

\begin{verbatim}
CAVITATION BUBBLE COLLAPSE:

    Time 1:           Time 2:           Time 3:
    
    ┌─────────┐      ┌───────┐         ┌─┐
    │         │      │       │         │●│ ← Extreme
    │    ○    │  →   │  ○    │   →     │ │   compression
    │         │      │       │         └─┘
    └─────────┘      └───────┘
    
    Bubble expands    Contracts         Collapse!
\end{verbatim}

At collapse point:
\begin{itemize}
    \item Extreme local pressure and temperature
    \item Membrane strongly deformed
    \item Possible Bulk channel opening
\end{itemize}

\textbf{Sonofusion (Taleyarkhan)} --- controversial but potentially EDC-explainable.

\subsection{Laser-Driven Focusing}

Femtosecond laser pulses can create coherent excitations:

\begin{verbatim}
LASER PULSE FOCUSING:

    Pulse 1 ───→ ╲
    Pulse 2 ───→  ╲
    Pulse 3 ───→   → TARGET
    Pulse 4 ───→  ╱
    Pulse 5 ───→ ╱
    
    Timing controls phase → focusing in time and space
\end{verbatim}

Multiple beams with controlled timing create interference at target.

%==============================================================================
\section{Theoretical Predictions}
%==============================================================================

\subsection{Testable Predictions}

If EDC is correct and coherent Bulk focusing works:

\begin{prediction}
\textbf{Prediction 1: Phase Sensitivity}
\begin{itemize}
    \item Fusion rate depends sensitively on relative phases
    \item Random phases: baseline rate
    \item Optimized phases: enhanced rate (potentially orders of magnitude)
\end{itemize}
\end{prediction}

\begin{prediction}
\textbf{Prediction 2: $N^2$ Scaling}
\begin{itemize}
    \item Doubling the number of coherent sources $\to$ 4$\times$ fusion rate
    \item Not 2$\times$ (which would indicate incoherent addition)
\end{itemize}
\end{prediction}

\begin{prediction}
\textbf{Prediction 3: Geometric Dependence}
\begin{itemize}
    \item Fusion rate depends on array geometry
    \item Optimal geometry focuses at target location
    \item Wrong geometry gives no enhancement
\end{itemize}
\end{prediction}

\begin{prediction}
\textbf{Prediction 4: Frequency Resonances}
\begin{itemize}
    \item Sharp peaks in fusion rate at specific frequencies
    \item Resonances have measurable width (Q-factor)
    \item Harmonic structure present
\end{itemize}
\end{prediction}

\begin{prediction}
\textbf{Prediction 5: Temperature Independence}
\begin{itemize}
    \item At resonance, fusion occurs below thermal threshold
    \item ``Cold fusion'' becomes possible at the right frequency
\end{itemize}
\end{prediction}

\subsection{Null Results}

If these predictions fail:
\begin{itemize}
    \item Coherent Bulk focusing doesn't work as proposed
    \item Membrane dynamics differ from model
    \item EDC requires modification
\end{itemize}

\textbf{This is a genuine scientific test.}

%==============================================================================
\section{Engineering Pathway}
%==============================================================================

\subsection{Phase 1: Proof of Concept}

\textbf{Goal:} Demonstrate any enhancement from coherent excitation

\begin{itemize}
    \item Small array (10--100 elements)
    \item Scanning frequency and phase
    \item Measure: neutron rate vs. random-phase baseline
\end{itemize}

\textbf{Success criterion:} Statistically significant enhancement with optimized phases

\subsection{Phase 2: Optimization}

\textbf{Goal:} Find optimal configuration

\begin{itemize}
    \item Systematic parameter search
    \item Machine learning for phase optimization
    \item Identify resonant frequencies
\end{itemize}

\textbf{Success criterion:} $>$10$\times$ enhancement over baseline

\subsection{Phase 3: Scaling}

\textbf{Goal:} Demonstrate net energy gain

\begin{itemize}
    \item Large arrays (1000+ elements)
    \item Optimized frequencies and phases
    \item Engineering for efficiency
\end{itemize}

\textbf{Success criterion:} $Q > 1$ (energy out $>$ energy in)

\subsection{Phase 4: Prototype Reactor}

\textbf{Goal:} Practical energy production

\begin{itemize}
    \item Continuous operation
    \item Heat extraction
    \item Grid integration
\end{itemize}

\textbf{Success criterion:} Reliable, economical power production

%==============================================================================
\section{Comparison with Current Approaches}
%==============================================================================

\begin{table}[H]
\centering
\begin{tabular}{@{}llll@{}}
\toprule
\textbf{Aspect} & \textbf{ITER (Tokamak)} & \textbf{NIF (Laser)} & \textbf{EDC Focusing} \\
\midrule
Temperature & 150 M K & 100 M K & Potentially room temp \\
Confinement & Magnetic & Inertial & None needed \\
Energy input & Enormous & Enormous & Modest (acoustic/EM) \\
Complexity & Extreme & Extreme & Moderate \\
Size & Building & Building & Tabletop possible \\
Fuel & D-T & D-T & D-D possible \\
Mechanism & Thermal & Thermal & Geometric \\
\bottomrule
\end{tabular}
\caption{Comparison of fusion approaches.}
\end{table}

\subsection{Why This Could Be Different}

Traditional fusion fails because it fights physics---trying to overcome barriers through brute force.

EDC fusion works \textbf{with} physics---finding the geometric path around barriers.

\begin{insight}
This is the difference between:
\begin{itemize}
    \item Climbing a mountain (thermal fusion)
    \item Taking a tunnel through the mountain (EDC fusion)
\end{itemize}
\end{insight}

%==============================================================================
\section{Societal Implications}
%==============================================================================

\subsection{If This Works}

Abundant, clean energy would transform human civilization:

\begin{itemize}
    \item \textbf{Climate:} Unlimited clean energy solves carbon emissions
    \item \textbf{Poverty:} Cheap energy enables development
    \item \textbf{Water:} Desalination becomes economical
    \item \textbf{Space:} Energy-intensive propulsion becomes feasible
    \item \textbf{Economy:} Energy scarcity ends as a constraint
\end{itemize}

\subsection{Ethical Commitment}

\begin{dedication}
This knowledge must not be monopolized.

We release this work under open license. We explicitly reject:
\begin{itemize}
    \item Patents on the fundamental mechanism
    \item Corporate ownership of the technique
    \item Restriction of access for profit
\end{itemize}

\textbf{Energy is a human right. The Sun belongs to everyone.}
\end{dedication}

\subsection{Call for Collaboration}

This theory needs experimental verification.

We invite:
\begin{itemize}
    \item \textbf{Physicists:} Test the predictions
    \item \textbf{Engineers:} Build the devices
    \item \textbf{Institutions:} Fund the research
    \item \textbf{Everyone:} Demand open access to results
\end{itemize}

%==============================================================================
\section{Conclusions}
%==============================================================================

We have proposed a new approach to nuclear fusion based on the 5D geometry of Elastic Diffusive Cosmology:

\begin{enumerate}
    \item \textbf{Quantum tunneling is Bulk transit} --- particles bypass barriers through higher dimensions
    \item \textbf{Bulk transit can be facilitated} --- by modifying local membrane conditions
    \item \textbf{Coherent focusing works in 5D} --- phased arrays can create focal points in the Bulk
    \item \textbf{At focal points, tunneling is enhanced} --- fusion becomes possible at lower temperatures
    \item \textbf{This is testable} --- specific predictions distinguish EDC from random effects
\end{enumerate}

The path forward is clear:
\begin{enumerate}
    \item Build phased arrays
    \item Scan for resonances
    \item Optimize configuration
    \item Scale to practical power
\end{enumerate}

\begin{keyresult}
If EDC is correct, we hold the key to unlimited clean energy.

\textbf{The stars run on fusion. Now, perhaps, so can we.}
\end{keyresult}

%==============================================================================
\section*{References}
%==============================================================================

\begin{enumerate}
    \item Grčman, I. (2026). ``Elastic Diffusive Cosmology - Part I.'' Zenodo. DOI: 10.5281/zenodo.18176174
    
    \item Grčman, I. (2026). ``Chemical Bonding from 5D Membrane Geometry.'' [This series, Paper I]
    
    \item Grčman, I. (2026). ``Jeans Mass from 5D Membrane Geometry.'' [This series, Paper II]
    
    \item Grčman, I. (2026). ``Nuclear Fusion from 5D Geometry.'' [This series, Paper III]
    
    \item Gamow, G. (1928). ``Zur Quantentheorie des Atomkernes.'' \textit{Z. Physik}, 51, 204.
    
    \item Taleyarkhan, R.P. et al. (2002). ``Evidence for Nuclear Emissions During Acoustic Cavitation.'' \textit{Science}, 295, 1868.
    
    \item Fleischmann, M. \& Pons, S. (1989). ``Electrochemically Induced Nuclear Fusion of Deuterium.'' \textit{J. Electroanal. Chem.}, 261, 301.
    
    \item Forsley, L.P. et al. (2020). ``A Review of Low Energy Nuclear Reactions.'' \textit{J. Condensed Matter Nucl. Sci.}, 33, 1.
\end{enumerate}

%==============================================================================
\appendix
\section{Detailed Array Design}
%==============================================================================

\subsection{Spherical Array Configuration}

For optimal 3D focusing, arrange transducers on a sphere:

\textbf{Geometry:}
\begin{itemize}
    \item Radius: $R = 10$ cm (adjustable)
    \item Elements: $N = 256$ (16 $\times$ 16 grid on sphere)
    \item Element size: $\sim$1 cm diameter
    \item Focal point: Center of sphere
\end{itemize}

\textbf{Phase calculation:}

For element $i$ at position $\vec{r}_i$, targeting focal point $\vec{r}_f$:
\begin{equation}
\phi_i = -\frac{2\pi f}{c}|\vec{r}_i - \vec{r}_f| + \phi_0
\end{equation}

\subsection{Electronics}

\begin{itemize}
    \item Function generators: $N$ channels, synchronized
    \item Phase resolution: $<1°$
    \item Frequency range: 1 kHz -- 100 MHz
    \item Amplitude control: 0--100\%
\end{itemize}

\subsection{Diagnostics}

\begin{itemize}
    \item Neutron detector: $^3$He proportional counter
    \item Gamma spectrometer: NaI or HPGe
    \item Calorimeter: Heat output measurement
    \item Oscilloscope: Verify phase coherence
\end{itemize}

%==============================================================================
\section{Safety Considerations}
%==============================================================================

\subsection{Radiation}

If fusion occurs, expect:
\begin{itemize}
    \item Neutrons (2.45 MeV from D-D, 14.1 MeV from D-T)
    \item Gamma rays
    \item Tritium (if D-T fuel)
\end{itemize}

\textbf{Shielding:}
\begin{itemize}
    \item Polyethylene/water for neutron moderation
    \item Lead for gamma
    \item Proper handling protocols for tritium
\end{itemize}

\subsection{Electrical}

High-voltage drivers for piezoelectric arrays:
\begin{itemize}
    \item Proper insulation
    \item Interlock systems
    \item Ground fault protection
\end{itemize}

%==============================================================================
\section{Estimated Costs}
%==============================================================================

\begin{table}[H]
\centering
\begin{tabular}{@{}ll@{}}
\toprule
\textbf{Item} & \textbf{Estimated Cost} \\
\midrule
Piezoelectric array (256 elements) & \$50,000 \\
Phase control electronics & \$100,000 \\
Vacuum system & \$30,000 \\
Neutron detector & \$20,000 \\
Gamma spectrometer & \$30,000 \\
Miscellaneous & \$20,000 \\
\midrule
\textbf{Total (basic setup)} & \textbf{$\sim$\$250,000} \\
\bottomrule
\end{tabular}
\caption{Estimated costs for basic experimental setup.}
\end{table}

This is orders of magnitude less than tokamak or laser fusion facilities.

%==============================================================================
\section{Open Source Commitment}
%==============================================================================

All designs, data, and results from experiments based on this work should be released publicly under open licenses (CC BY or CC0).

\textbf{Suggested protocol:}
\begin{enumerate}
    \item Pre-register experiments
    \item Publish all data (positive and negative)
    \item Share designs freely
    \item No patent applications on core mechanism
    \item Collaborate openly
\end{enumerate}

\vspace{1cm}
\hrule
\vspace{0.5cm}

\noindent\textit{Corresponding author: Igor Grčman}

\noindent\textit{Manuscript completed: January 2026}

\noindent\textit{This work is released under CC BY-NC 4.0 for the benefit of humanity.}

\vspace{1cm}

\begin{center}
\rule{0.8\textwidth}{0.4pt}

\textbf{Final Note}

\vspace{0.5cm}

This paper proposes what may seem like an extraordinary claim: that we can engineer nuclear fusion by manipulating higher-dimensional geometry.

Extraordinary claims require extraordinary evidence. We do not claim to have demonstrated this effect. We claim only to have:

\begin{enumerate}
    \item Identified a theoretical mechanism (EDC)
    \item Proposed a method to exploit it (coherent Bulk focusing)
    \item Outlined testable predictions
    \item Described experimental approaches
\end{enumerate}

The rest is up to experiment.

If we are right, the reward is unlimited clean energy.

If we are wrong, we will have learned something about the limits of EDC.

Either way, the experiment should be done.

\vspace{0.5cm}

\textit{\textbf{The universe has been running fusion for 13.8 billion years. It's time we learned the trick.}}

\rule{0.8\textwidth}{0.4pt}
\end{center}

\end{document}
